\documentclass[10pt]{article}
\usepackage[margin=20mm]{geometry}
\usepackage{amsmath,amssymb,mathtools,fontspec,unicode-math,xurl,hyperref}
\IfFontExistsTF{DejaVu Serif}{\setmainfont{DejaVu Serif}}{\setmainfont{Latin Modern Roman}}
\setmathfont{latinmodern-math.otf}
\setlength{\parindent}{0pt}\setlength{\parskip}{6pt}\sloppy
\title{Initial scientific review
Uniqueness of the 2D Euler equation on rough domains}
\author{Assigned initial-review executor}
\date{}
\begin{document}
\maketitle
Paper. Siddhant Agrawal and Andrea R. Nahmod, arXiv:2308.12926v1, 24 August 2023, 33 pages. Frozen identifier M25-0303; journal DOI mapping 10.4171/JEMS/1676. This report concerns the acquired preprint only.\par
\section*{Assessment}
No substantiated mathematical defect was established in the inspected argument. This is a bounded initial-review conclusion, not a certification of every imported theorem. All 33 pages were read as extracted text, including the appendix. Particular attention was given to the three conditions of Assumption 1.1, the construction of the angular change of variables, and the closure of the uniqueness estimate. The necessary bound source images and every report page were actually inspected in the same-chat continuation. The earlier image-return failure and incomplete delivery state are preserved in the verification history. The scientific decision, human review and independent adjudication remain pending.\par
\section*{Definitions, hypotheses and scope}
Theorem 1.2 concerns a bounded, simply connected Jordan domain and bounded initial vorticity. In conformal coordinates it requires a positive lower and upper bound for \(|\Psi_z|\), a uniform \(\dot H^{1/2}\) bound on the circular traces of \(\Psi_z\), and a uniform \(L^\infty\) bound for the Hilbert transform of \(|\Psi_z|^2\). None of these conditions was discarded in reviewing the proof.\par
The solution class is the Yudovich class of Definition 3.1, with the weak vorticity equation and orthogonality to gradient fields expressing incompressibility and the boundary condition. Simple connectivity is used in recovering the velocity from vorticity. The flow argument also uses conservation of the physical area measure, which becomes \(|\Psi_z(y)|^2\,dy\) on the disc. No uniqueness claim for arbitrary \(L^p\) vorticity or arbitrary \(C^{1,\alpha}\) domains with \(\alpha\leq1/2\) is inferred. The remarks about unbounded or multiply connected domains are extensions discussed by the authors, not separately verified theorems in this review.\par
\section*{Checks of the main proof}
\subsection*{Flow and boundary avoidance}
The Biot–Savart representation and its conformal transformation were followed in equations (12)–(20). Lemma 3.2 provides the bounded log-Lipschitz field \(\widetilde b\) and its vanishing normal component on the circle. Its boundary estimate is \(|\widetilde b\cdot\widehat r|\lesssim\phi(1-r)\), where \(\phi(q)=q\max\{-\log q,1\}\). Combined with the lower bound on \(|\Psi_z|\), this gives the Osgood control preventing trajectories from reaching the boundary in finite time. The local smoothing argument for measure preservation and the transport identity in Lemma 3.4 were read, with the external transport theorem treated as an imported dependency.\par
\subsection*{The change of variables}
The positive function \(G\), its circular average \(c(r)\) and \(L_\theta=G/c(r)\) give \(F(re^{i\theta})=re^{iL(r,\theta)}\), with \(F\) equal to the identity near the origin. Periodicity and strict monotonicity in the angular coordinate give a disc homeomorphism. Direct calculation gives \[\det DF=L_\theta=\frac{G}{c(r)}>0 .\] Together with the derivative bounds of Lemma 4.2 this controls both \(DF\) and \(DF^{-1}\) on the disc, as required by Lemma 4.3.\par
The nontrivial primitive estimate in Lemma 4.1 uses the real and imaginary parts of the holomorphic function \(\Psi_z\), the identity \(|\partial_\theta|a=b_\theta\), the quadratic identity (74), and the Hilbert-transform hypothesis. Thus the required bounded radial derivative is not being inferred merely from a bounded \(G\). The second-derivative bounds and Lemma 4.7 were traced to the \(\dot H^{1/2}\) and Poisson estimates.\par
Near the boundary, where \(G=|\Psi_z|^2\), the tangential coefficient has the exact cancellation \[\frac{|\Psi_z|^{-2}}{r}F_\theta
 =\frac{i e^{iL}}{c(r)} .\] The remaining radial coefficient has a difference bound involving the reciprocal distance to the boundary. On page 25 the proof splits according as the separation \(\delta\) is smaller or larger than that distance \(d\). In the first case \(\phi(d)\delta/d\lesssim\phi(\delta)\); in the second case the bounded coefficient and \(\phi(d)\leq\phi(\delta)\) suffice. These are the estimates that recover a log-Lipschitz bound despite the rough conformal derivative.\par
\subsection*{Energy closure and domain examples}
The three terms of the derivative of energy (63) were followed: an interior term, the transformed near-boundary term, and the difference between the two transported vorticities. The last uses the common initial vorticity, preservation of \(|\Psi_z|^2dy\), the reflected Biot–Savart kernel, and Fubini. The kernel comparison follows from \[|1-y\overline z|^2-|y-z|^2
 =(1-|y|^2)(1-|z|^2)\geq0 .\] The normalization to a probability measure before applying Jensen’s inequality is explicit on page 27. Small time ensures that the argument lies in the concavity interval of \(\phi\). The resulting inequality \(E'\lesssim\phi(E)\) and \(\int_{0+}dq/\phi(q)=+\infty\) yield uniqueness from \(E(0)=0\). The time step depends on the conserved vorticity bound and the domain, so iteration gives the global result.\par
For Proposition 1.3 the trace-composition argument, the strict threshold \(\alpha>1/2\), and the convex-domain tangent-angle variation bound were checked. Continuity of the logarithmic lift is used in Lemma 6.4; it is not an unrestricted lifting assertion for arbitrary measurable functions. Lemma 6.5 was read with its nonnegative, nonzero vorticity hypothesis; its lack-of-regularity claim is not being used as a general nonuniqueness result.\par
\section*{Independent checks and remaining limits}
An independently written SymPy program verifies the polar Jacobian identity, the reflected-kernel identity, a test-mode normalization of (74), and the divergent Osgood integral. These calculations check specific algebraic and limiting steps; they do not replace the analytic proof.\par
Two possible extraction-based concerns were rejected after consulting the version-matched arXiv mathematical HTML: Lemma 6.1 assumes \(|y'|\leq c\phi(y)\), and the coefficient in Lemma 4.5 is \(4g'/r-2g/r^2\). Neither is recorded as a paper defect. These two readings were also checked against the exact bound PDF images on pages 20 and 29 during the same-chat continuation.\par
The external weak-solution existence theorem, the full transport theorem and the conformal boundary regularity theorems were not independently reproved. A bounded title/erratum search and the author’s publication page did not reveal a separate correction, supplement or relevant code in the inspected results. Absence of code is not a mathematical defect, and no author code was run.\par
\section*{Supporting checks}
8403c2188c1dd5fb337941dc93b1f60c25286efb8d6c845d344708f75bb7b9ee\par
and the source-manifest SHA256 is\par
edc2b62630b6e6aad9514ffa626ae2db04ab007127e65194921bb1eb743f2283.\par
This is initial review, not independent acceptance. The necessary source images and all pages of this final report were actually inspected in the same-chat continuation. The previous incomplete state, original report and image-return failure remain preserved; the mathematical review and computations were not repeated.\par
9 S. Agrawal and A. R. Nahmod, Uniqueness of the 2D Euler equation on rough domains, arXiv:2308.12926v1 (2023).\par
\end{document}
